\chapter{Numerical Solutions to Partial Differential Equations}

\section{Classification of Second-Order PDEs}

\begin{definition}[Linear second-order PDE]
A linear second-order PDE in two variables has the form
\[
 Au_{xx}+Bu_{xy}+Cu_{yy}+Du_x+Eu_y+Fu=G,
\]
where the coefficients may depend on the independent variables.  It is classified by the discriminant $B^2-4AC$:
\begin{enumerate}[label=(\alph*),leftmargin=2.7em]
  \item parabolic if $B^2-4AC=0$;
  \item hyperbolic if $B^2-4AC>0$;
  \item elliptic if $B^2-4AC<0$.
\end{enumerate}
\end{definition}

\begin{example}
Classify the one-dimensional wave equation, the one-dimensional heat equation and the two-dimensional Laplace equation.
\end{example}

\begin{solution}[4.5cm]
For the wave equation
\[
 u_{tt}=c^2u_{xx},
\]
write $c^2u_{xx}-u_{tt}=0$.  Then $A=c^2$, $B=0$, $C=-1$, so
\[
 B^2-4AC=4c^2>0.
\]
It is hyperbolic.

For the heat equation
\[
 u_t=c^2u_{xx},
\]
the coefficient of the second derivative in time is zero, giving discriminant $0$.  It is parabolic.

For Laplace's equation
\[
 u_{xx}+u_{yy}=0,
\]
$A=C=1$ and $B=0$, so
\[
 B^2-4AC=-4<0.
\]
It is elliptic.
\end{solution}

\section{Finite Differences for Two-Variable Functions}

Let
\[
 x_i=x_0+ih,
 \qquad
 t_j=t_0+jk,
 \qquad
 u_{i,j}\approx u(x_i,t_j).
\]
Typical formulas are
\[
 u_x(x_i,t_j)
 =\frac{u_{i+1,j}-u_{i,j}}h+\bigO(h),
\]
\[
 u_t(x_i,t_j)
 =\frac{u_{i,j+1}-u_{i,j}}k+\bigO(k),
\]
\[
 u_x(x_i,t_j)
 =\frac{u_{i+1,j}-u_{i-1,j}}{2h}+\bigO(h^2),
\]
\[
 u_{xx}(x_i,t_j)
 =\frac{u_{i+1,j}-2u_{i,j}+u_{i-1,j}}{h^2}
 +\bigO(h^2).
\]
A finite-difference scheme replaces every derivative in the PDE by such algebraic expressions and then solves for the grid values.

\section{Parabolic Problems: The Heat Equation}

Consider
\[
 u_t=u_{xx},
 \qquad0<x<a,
 \qquad t>0,
\]
with an initial temperature distribution and boundary values.

\subsection{FTCS Explicit Scheme}

Using a forward difference in time and a central difference in space gives
\[
 \frac{u_{i,j+1}-u_{i,j}}k
 =\frac{u_{i+1,j}-2u_{i,j}+u_{i-1,j}}{h^2}.
\]
Define
\[
 r=\frac{k}{h^2}.
\]
Then
\[
 \boxed{
 u_{i,j+1}
 =(1-2r)u_{i,j}
 +r(u_{i+1,j}+u_{i-1,j})
 }.
\]
The method is explicit because all values at level $j+1$ are computed directly from the known level $j$.

\begin{example}
For
\[
 u_t=u_{xx},
 \qquad0\le x\le1,
\]
with
\[
 u(x,0)=\sin(\pi x),
 \qquad
 u(0,t)=u(1,t)=0,
\]
use FTCS with $h=0.1$ and $k=0.0005$ to estimate $u(0.1,0.001)$.
\end{example}

\begin{solution}[5.0cm]
Here
\[
 r=\frac{0.0005}{0.1^2}=0.05.
\]
At time $t=0$,
\[
 u_{1,0}=\sin(0.1\pi)=0.309017,
 \quad
 u_{2,0}=\sin(0.2\pi)=0.587785,
 \quad u_{0,0}=0.
\]
The first time step gives
\[
 u_{1,1}=0.9u_{1,0}+0.05(u_{2,0}+u_{0,0})
 =0.307505.
\]
After computing the neighbouring values at $t=0.0005$, the second step gives
\[
 u_{1,2}=0.3059995.
\]
Therefore
\[
 u(0.1,0.001)\approx0.3060.
\]
\end{solution}

\begin{theorem}[FTCS stability condition]
For the one-dimensional heat equation, the FTCS scheme is stable only when
\[
 0<r=\frac{k}{h^2}\le\frac12.
\]
\end{theorem}

Thus refining the spatial grid forces a much smaller time step.  If $h$ is halved, the stability limit requires $k$ to be reduced by a factor of four.

\subsection{Crank--Nicolson Implicit Method}

Crank--Nicolson averages the spatial central differences at time levels $j$ and $j+1$:
\[
 \frac{u_{i,j+1}-u_{i,j}}k
 =\frac12\left[
 \frac{u_{i+1,j+1}-2u_{i,j+1}+u_{i-1,j+1}}{h^2}
 +
 \frac{u_{i+1,j}-2u_{i,j}+u_{i-1,j}}{h^2}
 \right].
\]
With $r=k/h^2$, rearrangement gives
\[
 -ru_{i-1,j+1}+2(1+r)u_{i,j+1}-ru_{i+1,j+1}
 =ru_{i-1,j}+2(1-r)u_{i,j}+ru_{i+1,j}.
\]
At every new time level, the interior values are obtained from a tridiagonal linear system.  Crank--Nicolson is second-order accurate in both space and time and has no analogous restriction $r\le1/2$ for stability.

\begin{example}
For the same heat problem, use Crank--Nicolson with $h=0.2$ and $k=0.001$ to estimate $u(0.4,0.001)$.
\end{example}

\begin{solution}[6.0cm]
Here
\[
 r=\frac{0.001}{0.2^2}=0.025.
\]
There are four interior unknowns at the new time level.  They satisfy
\[
\begin{pmatrix}
2.05&-0.025&0&0\\
-0.025&2.05&-0.025&0\\
0&-0.025&2.05&-0.025\\
0&0&-0.025&2.05
\end{pmatrix}
\begin{pmatrix}
u_{1,1}\\u_{2,1}\\u_{3,1}\\u_{4,1}
\end{pmatrix}
=
\begin{pmatrix}
1.169958\\1.893031\\1.893031\\1.169958
\end{pmatrix}.
\]
Solving gives
\[
 (u_{1,1},u_{2,1},u_{3,1},u_{4,1})
 \approx(0.582199,0.942018,0.942018,0.582199).
\]
Since $x=0.4$ corresponds to $i=2$,
\[
 u(0.4,0.001)\approx0.9420.
\]
\end{solution}

For this problem, the exact solution is
\[
 u(x,t)=e^{-\pi^2t}\sin(\pi x).
\]
This allows explicit and implicit schemes to be compared with the true solution.  FTCS can be very accurate when its stability condition is respected, but it may become unstable when $k/h^2$ is too large.  Crank--Nicolson remains stable and usually gives better accuracy at larger time steps.

\section{Elliptic Problems: Laplace's Equation}

For
\[
 u_{xx}+u_{yy}=0,
\]
use central differences on a rectangular grid with spacings $h$ and $k$:
\[
 \frac{u_{i+1,j}-2u_{i,j}+u_{i-1,j}}{h^2}
 +
 \frac{u_{i,j+1}-2u_{i,j}+u_{i,j-1}}{k^2}=0.
\]
For a square grid $h=k$, this simplifies to the five-point formula
\[
 \boxed{
 u_{i,j}=\frac14
 \left(u_{i+1,j}+u_{i-1,j}+u_{i,j+1}+u_{i,j-1}\right)
 }.
\]
Thus each interior temperature is the average of its four nearest neighbours.  All interior equations form a sparse linear system, commonly solved by Jacobi, Gauss--Seidel or SOR.

\begin{example}
A square plate of side $12$ cm has temperatures
\[
 \text{left}=	ext{right}=	ext{bottom}=100^\circ\mathrm C,
 \qquad
 \text{top}=0^\circ\mathrm C.
\]
Using a mesh size of $4$ cm, find the temperatures at
\[
 A(4,4),\quad B(8,4),\quad C(8,8),\quad D(4,8).
\]
Gauss--Seidel starts from
\[
 u_A=u_B=80,
 \qquad u_C=u_D=50,
\]
and stops when every change is below $10^{-3}$.
\end{example}

\begin{solution}[7.0cm]
The five-point equations are
\[
 u_A=\frac{100+u_B+u_D+100}{4},
\]
\[
 u_B=\frac{u_A+100+u_C+100}{4},
\]
\[
 u_C=\frac{u_D+100+0+u_B}{4},
\]
\[
 u_D=\frac{100+u_C+0+u_A}{4}.
\]
By symmetry, $u_A=u_B$ and $u_C=u_D$.  Let these common values be $a$ and $c$.  Then
\[
 a=50+\frac{a+c}{4},
 \qquad
 c=25+\frac{a+c}{4}.
\]
Equivalently,
\[
 3a-c=200,
 \qquad
 -a+3c=100.
\]
Solving gives
\[
 a=87.5,
 \qquad c=62.5.
\]
Gauss--Seidel converges to the same values:
\[
 u_A=u_B=87.5^\circ\mathrm C,
 \qquad
 u_C=u_D=62.5^\circ\mathrm C.
\]
\end{solution}

\section{Hyperbolic Problems: The Wave Equation}

Consider
\[
 u_{tt}=u_{xx},
\]
with fixed-end boundary conditions and initial displacement and velocity:
\[
 u(0,t)=u(a,t)=0,
\]
\[
 u(x,0)=f(x),
 \qquad
 u_t(x,0)=g(x).
\]
Using central differences in both time and space gives
\[
 \frac{u_{i,j+1}-2u_{i,j}+u_{i,j-1}}{k^2}
 =
 \frac{u_{i+1,j}-2u_{i,j}+u_{i-1,j}}{h^2}.
\]
Let
\[
 \rho=\left(\frac{k}{h}\right)^2.
\]
Then the CTCS scheme is
\[
 \boxed{
 u_{i,j+1}
 =\rho u_{i+1,j}+2(1-\rho)u_{i,j}
 +\rho u_{i-1,j}-u_{i,j-1}
 }.
\]
The stability condition is
\[
 0\le\rho\le1,
 \qquad\text{equivalently }k\le h.
\]

The initial velocity requires a special first time level.  A central approximation at $t=0$ gives
\[
 \frac{u_{i,1}-u_{i,-1}}{2k}=g_i.
\]
Eliminating the fictitious value $u_{i,-1}$ from the CTCS equation yields
\[
 u_{i,1}
 =\frac\rho2(f_{i+1}+f_{i-1})
 +(1-\rho)f_i+kg_i.
\]

\begin{example}
A string with fixed ends satisfies
\[
 u_{tt}=u_{xx},
 \qquad0\le x\le1,
\]
with
\[
 u(x,0)=0,
 \qquad
 u_t(x,0)=\sin(\pi x),
\]
and $u(0,t)=u(1,t)=0$.  Use CTCS with $h=k=0.2$ to find the displacement at $t=0.4$ for
$x=0.2,0.4,0.6,0.8$.
\end{example}

\begin{solution}[6.2cm]
Since $h=k$, we have $\rho=1$.  The initial displacement is $f_i=0$, so the first time level is
\[
 u_{i,1}=kg_i=0.2\sin(\pi x_i).
\]
Thus at $t=0.2$,
\[
 (u_{1,1},u_{2,1},u_{3,1},u_{4,1})
 \approx(0.11756,0.19021,0.19021,0.11756).
\]
For the next level, because $\rho=1$ and $u_{i,0}=0$,
\[
 u_{i,2}=u_{i+1,1}+u_{i-1,1}.
\]
Using the fixed boundary values $u_{0,1}=u_{5,1}=0$, we obtain
\[
\begin{aligned}
 u_{1,2}&=0.19021,\\
 u_{2,2}&=0.11756+0.19021=0.30777,\\
 u_{3,2}&=0.19021+0.11756=0.30777,\\
 u_{4,2}&=0.19021.
\end{aligned}
\]
Therefore at $t=0.4$ the displacements are approximately
\[
 0.1902,
 \quad0.3078,
 \quad0.3078,
 \quad0.1902.
\]
\end{solution}
